\textbf{Proof.}
Fix \(\eta\in\mathfrak C_D(\mathbf M)\), and put
\[
E=\mathcal G_\eta[V],
\qquad
U_\eta=\bigcup_{s\in S_\eta}\operatorname{pa}_{\mathcal G_\eta}(s),
\qquad
A=\operatorname{an}_E(U_\eta).
\]
By \(\operatorname{ass{:}selection\text{-}parents}\), \(U_\eta\subseteq V\).  Hence \(A\subseteq V\), and idempotence of ancestor closure gives \(\operatorname{an}_E(A)=A\).  The restriction \(E\) of the finite DAG \(\mathcal G_\eta\) is acyclic.  If \(\{i,j\}\in Q_\eta\), then \(i,j\in U_\eta\subseteq A\), so \(Q_\eta\subseteq\binom A2\).

By the observational formula in \(\texttt{lem:dai-two-section-checker-characterization}\), the observational projection of \(\eta\) has adjacency set \(B(E,A)\cup Q_\eta\).  Membership in \(\mathfrak C_D(\mathbf M)\) entails observational Markov equivalence to \(\mathcal M^{(0)}\), and Markov-equivalent MAGs have the same adjacencies.  Therefore
\[
B(E,A)\cup Q_\eta=S_0(\mathbf M).
\]
It follows that
\[
B(E,A)\subseteq S_0(\mathbf M),
\qquad
Q_{\mathbf M}(E,A)
=S_0(\mathbf M)\setminus B(E,A)
\subseteq Q_\eta
\subseteq\binom A2.
\]

Fix \(k\in\mathcal K_+\), write \(T_k=I_\eta^{(k)}\), and let \(F_k=\operatorname{de}_E(T_k)\).  The indicator formula in \(\texttt{lem:dai-two-section-checker-characterization}\) and compatibility give
\[
Z_k(\mathbf M)=T_k\cup(F_k\cap A).
\]
Because \(F_k\cap A\subseteq A\), taking the difference with \(A\) gives
\[
O_k(\mathbf M,A)=Z_k(\mathbf M)\setminus A=T_k\setminus A.
\]
Since \(T_k\subseteq F_k\), intersecting the indicator identity with \(A\) gives
\[
R_k(\mathbf M,A)=Z_k(\mathbf M)\cap A=F_k\cap A.
\]
If \(r\in R_k(\mathbf M,A)\) and \(a\in\operatorname{de}_E(r)\cap A\), transitivity of descendants gives \(a\in F_k\cap A=R_k(\mathbf M,A)\).  Thus \(R_k(\mathbf M,A)\) is descendant-closed in \(E[A]\).

Write \(R_k=R_k(\mathbf M,A)\), \(O_k=O_k(\mathbf M,A)\), and \(T_k^*=T^*_{k,\mathbf M}(E,A)\).  Since \(A\) is ancestor-closed, no vertex outside \(A\) has a descendant in \(A\), and so
\[
\operatorname{de}_E(O_k)\cap A=\varnothing.
\]
The finite-DAG descent argument from the preceding lemma gives
\[
\operatorname{de}_E(\operatorname{Min}_E(R_k))\cap A=R_k.
\]
Every \(m\in\operatorname{Min}_E(R_k)\) lies in \(R_k=\operatorname{de}_E(T_k)\cap A\).  No member of \(T_k\setminus A=O_k\) can reach \(A\), so \(m\) descends from some vertex in \(T_k\cap A\).  Hence
\[
\operatorname{de}_E(\operatorname{Min}_E(R_k))
\subseteq\operatorname{de}_E(T_k\cap A).
\]
Conversely, \(T_k\cap A\subseteq R_k\), and every element of the finite descendant-closed set \(R_k\) descends from a member of \(\operatorname{Min}_E(R_k)\).  Thus the reverse inclusion holds.  Therefore
\[
\begin{aligned}
\operatorname{de}_E(T_k^*)
&=\operatorname{de}_E(O_k\cup\operatorname{Min}_E(R_k))\\
&=\operatorname{de}_E(T_k\setminus A)\cup\operatorname{de}_E(T_k\cap A)\\
&=\operatorname{de}_E(T_k)=F_k.
\end{aligned}
\]

Let \(Q^*=Q_{\mathbf M}(E,A)\).  Retain \(E\), add one binary childless selection node with parent set \(\{i,j\}\) for each \(\{i,j\}\in Q^*\), and add one binary childless selection node with parent set \(\{i\}\) for every \(Q^*\)-isolated \(i\in A\).  The union of reconstructed parent sets is \(A\), its two-section is \(Q^*\), and \(B(E,A)\cup Q^*=S_0(\mathbf M)\).  Replace each nonobservational target \(T_k\) by \(T_k^*\).  The observational projection is unchanged.  The equality \(\operatorname{de}_E(T_k^*)=F_k\) preserves every intervention-induced path adjacency and every nonindicator endpoint, while
\[
T_k^*\cup(\operatorname{de}_E(T_k^*)\cap A)
=O_k\cup R_k
=Z_k(\mathbf M)
\]
preserves every indicator adjacency.  The checker characterization therefore gives equality of the fully oriented projected MAGs in every regime.

The reconstructed graph is finite and acyclic because the new selection nodes are childless.  It is causally sufficient on \(V\); every selection node is binary, childless, and observed-parented; the selected stratum, pre-treatment timing, and empty observational target are retained; and soft-intervention modularity holds for the new ordered target family.  Choose every binary selection kernel strictly between zero and one on every parent configuration.  Finiteness then gives positive probability to the joint selected stratum.  Thus \(\bar\eta\in\mathfrak C_D(\mathbf M)\).

Conversely, a normalized target-residual certificate already derives its unique target family from \((\mathbf M,E,A)\).  The same pair-node and isolated-singleton construction has selection-parent union \(A\) and two-section \(Q_{\mathbf M}(E,A)\), and uses at most
\[
\lvert Q_{\mathbf M}(E,A)\rvert+\lvert A\rvert
\le \binom D2+D
\]
selection nodes.  When \(K=0\), \(\mathcal K_+\) is empty, so no target argument or replacement occurs.  No overlap, rank, inverse, division, or limiting condition is used. \(\square\)